
\newcounter{phan}[section]
\newcommand{\phan}[2][\mycolor]{%
	\refstepcounter{phan}
	\par\noindent\begin{tikzpicture}[declare function={d=0.2cm;r=0.6cm;}]
		\path (0,0) node[font=\fontfamily{qag}\selectfont\bfseries\color{white},inner sep=2pt,outer sep =2pt](PHAN){\Alph{phan}.\,#2};
		\path (PHAN.north east) coordinate (A)
		      ([xshift=r]PHAN.south east) coordinate (B);
		\foreach \a/\d in {0/80,3/65,6/50,9/35,12/20}{
		\begin{scope}[transform canvas={xshift=\a pt}]
			\path [fill=#1!\d] ([xshift=1pt]A)--([xshift=1pt]B)--([xshift=3pt]B)--([xshift=3pt]A)--cycle;
		\end{scope}
		}
		\path[left color=#1!90,right color=#1!10] (PHAN.south west) rectangle++({\linewidth-d},2pt);
		\path[fill=#1!80] (PHAN.north west)--(PHAN.south west)-- ([xshift=r]PHAN.south east)--(PHAN.north east)--cycle;
		\path (0,0) node[font=\fontfamily{qag}\selectfont\bfseries\color{white},inner sep=2pt,outer sep =2pt](PHAN){\Alph{phan}.\,#2};
		\path[fill=#1!80!black] (PHAN.south west)--(PHAN.north west)--([turn]90:d)--([turn]90:d)|-(PHAN.south west);
	\end{tikzpicture}%	
	\vspace{-0.5cm}
}

%%%=========================Dạng 1====================%%%
\newcounter{dang}[section]
\NewTColorBox{dang}{m}{
	enhanced jigsaw,
	before skip = 6pt,after skip =-6pt,
	code={%
		\refstepcounter{dang}
		\addcontentsline{toc}{subsection}{\faBook\ Dạng \thedang. #1}
		\setcounter{subsubsection}{0}
		\setcounter{phan}{0}
		},
	empty,
	title={\faBook\ Dạng \thedang. #1},
	fonttitle= \color{\mycolor}\bfseries\large\fontfamily{phv}\selectfont,
	breakable,
	arc is angular,
	boxsep =6pt, left=3pt, top=6pt,right=1pt,
	overlay unbroken and first={%
		 \path([yshift=-1mm]title.south west) rectangle (frame.south east); 
		 \path[fill=\mycolor!9,draw=\mycolor!50!black,rounded corners =6pt](title.north west) rectangle (title.south east);
	},
	overlay middle and last={%
		\path([yshift=-1mm]frame.north west) rectangle (frame.south east); 
	}
}
%%%==================end dang======================%%%

%%% Tùy chọn 1: màu có thể bỏ trống
%%%\tùy chọn 2: font chữ có thể bỏ trống
%%%\tùy chọn 3: nhãn có thể bỏ trống
\DeclareFixedFont{\ugqfont}{T5}{ugq}{b}{n}{13pt}
\newcommand{\qagfont}{\bfseries\fontfamily{qag}\fontsize{13pt}{6pt}\selectfont}
\newcounter{noibat}[subsubsection]
\renewcommand{\thenoibat}{\alph{noibat}}
\NewDocumentCommand{\Noibat}{O{}O{}O{}O{}O{0}m}{%
	% Chỉ tăng số đếm khi sử dụng tùy chọn \thenoibat hoặc khi #3 trống
	\ifstrequal{#3}{\thenoibat}{\refstepcounter{noibat}}{%
		\ifblank{#3}{\refstepcounter{noibat}}{}%
	}%
	\ifblank{#1}{\def\tuychonone{\mycolor}}{\def\tuychonone{#1}}%
	\ifblank{#2}{\def\tuychontwo{\qagfont}}{\def\tuychontwo{#2}}%
	\ifblank{#3}{\def\tuychonthree{\thenoibat)}}{\def\tuychonthree{#3}}%
	\ifblank{#4}{\def\tuychonfour{0}}{\def\tuychonfour{#4}}%
	\par\noindent\hspace{\tuychonfour cm}%
	{\color{\tuychonone}{\large\tuychontwo\expandafter\tuychonthree\,}{\tuychontwo #6}}%
	\vspace*{#5 cm}
}
\newcommand{\tieumuc}[1]{\noindent\Noibat[\maunhan][\qagfont][\faBookmark]{#1}}

\newcommand{\Nhanmanh}[2][\faComments]{\Noibat[\maunhan][\qagfont][#1]{#2}}


\newcommand\khungct[2][\mycolor]
{\noindent
	\begin{tikzpicture}[baseline=(char.base)]%
		\tikzstyle{khung} = [rectangle,draw=#1!30!black,fill=#1!10,inner sep=4pt,rounded corners]%
		\node [khung] (char){%
			$
			\; #2.\;
			$
		};%
\end{tikzpicture}}

\newcommand{\hopcttoan}[2][\maunhan]{%
	\tcboxmath[arc=3pt,arc is angular,colback=#1!10,boxrule=0.65pt,colframe=#1,boxsep=2pt,fontupper=\color{#1!50!black}]{#2}}

\NewDocumentCommand{\boxct}{O{\maunhan}O{3pt}O{}m}{
	\tcbox[%
	enhanced,
	ignore nobreak,
	arc=#2,
	colback=#1!15,
	boxrule=0.65pt,
	colframe=#1, 
	boxsep=3pt,
	left=2pt,right=2pt,
	tikznode upper,
	fontupper = #3,
	]{#4}
}


%\newcounter{tomtat}[section]
%\newtcolorbox{boxtomtat}{%
%	enhanced jigsaw,
%	breakable,
%	empty,
%	colback=\mycolor!5,
%	left=2mm,
%	right=1mm,
%	top=3mm,
%	bottom=3mm,
%	after skip=6pt,
%	before skip=6pt,
%	left skip=0.65cm,
%	toprule=0pt,
%	bottomrule=0pt,
%	rightrule=1.5pt,
%	leftrule=1.5pt,
%	overlay unbroken and first ={
%	\path[fill=\mycolor!5,draw=\mycolor!50!black,rounded corners =5pt] (frame.north west) rectangle (frame.south east);
%	},
%	overlay middle and last={
%		\draw[fill=\mycolor!5,draw=\mycolor!50!black,rounded corners =5pt] (frame.north west) rectangle (frame.south east);
%	}
%%	colframe=\mycolor,
%%	drop lifted shadow=\maudam!40!black
%	}
%	\newenvironment{tomtat}{\begin{boxtomtat}}{\end{boxtomtat}}

\newtcolorbox{ppgbox}[1][\mycolor]{
enhanced jigsaw,
left=2mm,
right=1mm,
top=3mm,
bottom=3mm,
after skip=6pt,
before skip=6pt,
arc=0pt,
breakable,
colbacktitle=#1!35,
empty,opacityback=0,%% Tách nền
boxed title style ={empty,sharp corners,boxrule =0pt,rounded corners=northwest,arc = 6pt,rounded corners=northeast},
title={Phương pháp giải},
fonttitle=\bfseries\fontfamily{phv}\selectfont, 
halign title=center,coltitle=#1!30!black,
attach boxed title to top center={xshift=0mm,yshift=0mm},
overlay unbroken and first={
	\path (frame.north west) coordinate (FNW)
	(frame.north east) coordinate (FNE)
	(frame.south east) coordinate (FSE)
	(frame.south west) coordinate (FSW)
	;
	\path [fill =#1!10,draw=#1] (title.south west)
	--([yshift=-6pt]title.north west)to[in=180,out=90]([xshift=6pt]title.north west)
	--([xshift=-6pt]title.north east)to[in=90,out=0]([yshift=-6pt]title.north east)--(title.south east)--cycle;
	\path[fill=#1,draw=#1,line width=1pt,rounded corners] ([xshift=8pt]FNW) --(title.south west)circle (2pt)
	(title.south east)circle (2pt)--([xshift=-8pt]FNE);
	\path[draw=#1,dashed]
	([xshift=-8pt]FNE)--([yshift=-5pt]FNE) --([yshift=5pt]FSE)
	--([xshift=-8pt]FSE) --([xshift=8pt]FSW)
	--([yshift=5pt]FSW) --([yshift=-5pt]FNW)
	--([xshift=8pt]FNW)
	;
	\path ($(title.west)!0.5!(title.east)$) node[font=\color{#1!30!black}\bfseries\fontfamily{phv}\selectfont]{Phương pháp giải};
	\path ($(title.east)+(0,0)$) node[anchor = west]{\includegraphics[height=1cm]{Images/Icons/PPG.PNG}};
	\path ($(title.west)+(0.2,0)$) node[anchor = east]{\includegraphics[height=0.8cm]{Images/Icons/PEN.PNG}};
},
overlay middle and last={
	\path (frame.north west) coordinate (FNW)
	(frame.north east) coordinate (FNE)
	(frame.south east) coordinate (FSE)
	(frame.south west) coordinate (FSW);
	\path[draw=#1,dashed]([xshift=8pt,yshift=-5pt]FNW)--([xshift=-8pt,yshift=-5pt]FNE)
	--([yshift=-10pt]FNE) --([yshift=5pt]FSE)
	--([xshift=-8pt]FSE) --([xshift=8pt]FSW)
	--([yshift=5pt]FSW) --([yshift=-10pt]FNW)
	--([xshift=8pt,yshift=-5pt]FNW)
	;
}
}
\newenvironment{pp}[1][\mycolor]{\begin{ppgbox}[#1]}{\end{ppgbox}}

\newenvironment{cacbuoc}[1][\maunhan]{
	\begin{enumerate}[label= \color{#1}\bfseries\fontfamily{qag}\selectfont{\faArrowCircleORight\;Bước \arabic*:},itemsep=3pt,wide=0cm,leftmargin=0cm,topsep=3pt]
	}{\end{enumerate}}

%%%================Các môi trường==============%%%
%%%================Môi trường Câu hỏi khởi động===================%%%
\NewEnviron{cauhoikd}{
	\begin{center}
		\begin{tikzpicture}
			\node[cloud,cloud ignores aspect=true,cloud puffs=35,cloud puff arc=70,draw=\mauphu!65,fill=white,
			text width=0.5\linewidth,font=\sffamily,drop shadow={opacity=0.25,color=\maudam}]
			(nA){\BODY};
			\begin{scope}[rotate=5]
				\draw[\mauphu,fill=white] ($(nA.south west)+(-2,0.5)$) coordinate (A) circle (0.6)
				($(A)+(240:0.375)$) arc (240:310:0.35)
				($(A)+(30:0.2)$) circle (2pt)
				($(A)+(150:0.2)$) circle (2pt);
				\draw[thin,\mauphu!65] (A) circle (0.675);
				\draw[\mauphu!85,line cap=round,thick,shift={(A)}] (15:0.675) arc (15:165:0.675)
				(-15:0.675) arc (-15:-165:0.675);
				\fill[\mauphu] ($(A)+(150:0.185)$) circle (1pt)
				($(A)+(30:0.185)$) circle (1pt);
				\path ($(A)+(60:0.7)+(-45:1pt)$) node[rotate=5,font=\bfseries\LARGE\color{gray}]{?};
				\path ($(A)+(60:0.7)$) node[rotate=5,font=\bfseries\huge\color{\mycolor!50!orange}]{?};
			\end{scope}
		\end{tikzpicture}
\end{center}}
%%%================Môi trường tri thức===================%%%
\NewEnviron{trithuc}{
	\begin{center}
		\begin{tikzpicture}
			\node[decorate,decoration={snake},segment length=15pt,draw=\maunhan!65,thick,
			fill=white,text width=0.8\linewidth,font=\sffamily,drop shadow={opacity=0.25,color=\maunhan}]
			(nA){\BODY};
			\begin{scope}[rotate=5]
				\draw[\maunhan,fill=white] ($(nA.north west)+(-1,-0.25)$) coordinate (A) circle (0.6)
				($(A)+(240:0.375)$) arc (240:310:0.35)
				($(A)+(30:0.2)$) circle (2pt)
				($(A)+(150:0.2)$) circle (2pt);
				\draw[thin,\maunhan!65] (A) circle (0.675);
				\draw[\maunhan!85,line cap=round,thick,shift={(A)}] (15:0.675) arc (15:165:0.675)
				(-15:0.675) arc (-15:-165:0.675);
				\fill[\maunhan] ($(A)+(150:0.185)$) circle (1pt)
				($(A)+(30:0.185)$) circle (1pt);
				\path ($(A)+(60:0.7)+(-45:1pt)$) node[rotate=-5,font=\bfseries\Huge\color{gray}]{\faLightbulbO};
				\path ($(A)+(60:0.7)$) node[rotate=-5,font=\bfseries\Huge\color{\maunhan!90!pink}]{\faLightbulbO};
			\end{scope}
		\end{tikzpicture}
\end{center}}
%%%============lệnh tìm hiểu thêm============%%%
\newcommand{\timhieuthem}{%
	\par
	\noindent
	\hspace*{-0.65cm}
	\begin{tikzpicture}[declare function={rt=0.65;d=sqrt(2)*rt;}]
		\path (90:rt) coordinate (A)
		(180:rt) coordinate (B)
		(270:rt) coordinate (C)
		(0:rt) coordinate (D)
		([shift={(0.15,0)}]0:rt) coordinate (E)
		($ (E)+(135:rt) $) coordinate (F)
		($ (E)+(-135:rt) $) coordinate (G)
		($ (A)!0.5!(F) $) coordinate (AF)
		($ (D)!0.5!(E) $) coordinate (DE)
		($ (C)!0.5!(G) $) coordinate (CG)
		;
		\path[transform canvas={shift={(-45:2.5pt)}},left color=\mycolor!50,right color=\mycolor!5]
		(F)--(E)--(G)--+(1.25*\linewidth,0)--([turn]90:d)--cycle;
		\path[outer color=\mycolor!65!black,inner color=\mycolor!80,rounded corners=2pt] (A)--(B)--(C)--(D)--cycle;
		\draw[line width=0.65pt,\mycolor,line cap=round,rounded corners=1.5pt]
		(AF)--(DE)--(CG);
		\draw[transform canvas={xscale=-1},line width=0.65pt,\mycolor,line cap=round,rounded corners=1.5pt] (AF)--(DE)--(CG);
		\path[left color=\mycolor!90!black,right color=\mycolor!50] (F)--(E)--(G)--+(1.25*\linewidth,0)--([turn]90:d)--cycle;
		\path (-0.7pt,-0.3pt) node[inner sep=0pt,font=\Large\itshape\bfseries\color{gray},xscale=-1]{\faSearchPlus};
		\path (-1pt,0) node[inner sep=0pt,font=\Large\itshape\bfseries\color{white},xscale=-1]{\faSearchPlus};
		\path ([shift={(0.11,-0.01)}]E) node[align=left,anchor=west,font=\large\sffamily\bfseries\color{gray}]{TÌM HIỂU THÊM};
		\path ([shift={(0.1,0)}]E) node[align=left,anchor=west,font=\large\sffamily\bfseries\color{white}]{TÌM HIỂU THÊM};
	\end{tikzpicture}%
}
%%%==============Môi trường Bài tập tương tự=====================%%%
\newtcolorbox{bttuongtu}{
	enhanced,breakable,
	before skip=6pt,after skip=6pt,
	boxsep=9pt,
	left=0.75cm,
	right=0.15cm,
	bottom=0pt,
	frame empty,
	frame hidden,
	arc=5pt,
	colback=\mycolor!7,
	sharp corners,
	rounded corners=south,
	colframe=\mycolor!30!\maunhan!45,
	boxrule=0pt,
	borderline north={9pt}{0pt}{\mycolor!40},
	drop fuzzy shadow=gray!75,
	overlay unbroken and first ={
		\fill[\mycolor!80!black,drop shadow={opacity=0.25,shadow xshift=1pt, shadow yshift=-1.5pt,shadow scale=1.05}] (frame.north west) coordinate (A)--+(1.5,0) coordinate (Bt)--($(A)+(0,-1.25)$)coordinate (Ct)--cycle;
		\path ($(A)+(0.5,-0.4)$) coordinate (B) node[font=\LARGE\color{gray!65}]{\faEdit};
		\path ($(B)+(-1pt,1pt)$) node[font=\LARGE\color{yellow!20}]{\faEdit};
		\path (Bt)--(Ct) node[below=-2.5pt,sloped,pos=0.5,text=\mycolor!75!\maunhan]{\scriptsize\sffamily bài tập tương tự};
	}
}


%%%============môi trường định nghĩa============%%%
\newcounter{dn}
\newenvironment{dn}[1][]{\setcounter{dn}{\arabic{subsubsection}}\refstepcounter{subsubsection}\refstepcounter{dn}\ifblank{#1}{{\sffamily\bfseries\color{\mycolor}\par\noindent\thedn) Định nghĩa.\,}}%
	{{\sffamily\bfseries\color{\mycolor}\par\noindent\thedn) Định nghĩa (#1).\,}}}{}
%%%============môi trường định lý============%%%
\newcounter{dl}
\newenvironment{dl}[1][]{\setcounter{dl}{\arabic{subsubsection}}\refstepcounter{subsubsection}\refstepcounter{dl}\ifblank{#1}{{\sffamily\bfseries\color{\mycolor}\par\noindent\thedl) Định lý.\,}}%
	{{\sffamily\bfseries\color{\mycolor}\par\noindent\thedl) Định lý (#1).\,}}\itshape}{}
%%%============môi trường hq============%%%
\newcounter{hq}
\newenvironment{hq}[1][]{\setcounter{hq}{\arabic{subsubsection}}\refstepcounter{subsubsection}\refstepcounter{hq}\ifblank{#1}{{\sffamily\bfseries\color{\mycolor}\par\noindent\thehq) Hệ quả.\,}}%
	{{\sffamily\bfseries\color{\mycolor}\par\noindent\thehq) Hệ quả (#1).\,}}\itshape}{}
%%%============môi trường tc============%%%
\newcounter{tc}
\newenvironment{tc}[1][]{\setcounter{tc}{\arabic{subsubsection}}\refstepcounter{subsubsection}\refstepcounter{tc}\ifblank{#1}{{\sffamily\bfseries\color{\mycolor}\par\noindent\thetc) Tính chất.\,}}%
	{{\sffamily\bfseries\color{\mycolor}\par\noindent\thetc) Tính chất (#1).\,}}\itshape}{}
%%%==========Môi trường hoạt động===============%%%
\tikzset{laphuong/.pic={
		\path[pic actions,line join=round]
		(0,0)--++(27.5:r)--++(152.5:r)--+(-152.5:r)--cycle
		(0,0)--++(0,-r)--++(152.5:r)--+(0,r)--cycle
		(0,0)--++(0,-r)--++(27.5:r)--+(0,r)--cycle;
}}
\newcounter{hd}[section]
\newtcolorbox{hd}{before skip=6pt,after skip=6pt,
	left skip=-0.35cm,
	boxsep=1pt,
	enhanced,
	top=0pt,
	frame empty,
	frame hidden,
	opacityback=0,
	code={\refstepcounter{hd}},
	before upper={
		\begin{tikzpicture}[declare function={r=0.165;},baseline=(char.base)]
			\path (30:1.1*r) pic[fill=orange,draw=white]{laphuong};
			\path (150:1.1*r) pic[fill=red!75!blue,draw=white]{laphuong};
			\path (270:1.075*r) pic[fill=green!75!blue,draw=white]{laphuong};
			\fill[\mycolor] (60:{2.5*r})--+(4*r,0)--(7.5*r,0)--($(-60:{2.5*r})+(4*r,0)$)--+(-4*r,0) arc(-60:60:{2.5*r});
			\draw[ultra thick,\mycolor](0,0) circle (2.5*r) (8*r,0);
			\path ($({4*r},0)+(0.5pt,-0.65pt)$) node[font=\bfseries\sffamily\color{gray!65}\large] (char){\thehd};
			\path ({4*r},0) node[font=\bfseries\sffamily\color{white}\large] (char){\thehd};
	\end{tikzpicture}}
}

%%%===============Môi trường kiến thức trong tâm====================%%%

\newtcolorbox{kttrongtam}{
	enhanced,breakable,
	before skip=6pt,after skip=6pt,
	boxsep=9pt,
	left=1cm,
	right=0.25cm,
	bottom=0pt,
	top=0.35cm,
	frame empty,
	frame hidden,
	arc=7pt,
	colback=\mycolor!7,
	sharp corners,
	rounded corners=south,
	colframe=\mycolor!35,
	boxrule=0pt,
	borderline north={9pt}{0pt}{\mycolor!65},
	drop fuzzy shadow=gray!75,
	overlay unbroken and first ={
		\fill[\mycolor,drop shadow={opacity=0.25,shadow xshift=1pt, shadow yshift=-1.5pt,shadow scale=1.05}] (frame.north west) coordinate (A)--+(2,0)--($(A)+(0,-1.25)$)--cycle;
		\path ($(A)+(0.65,-0.15)$) coordinate (B) node{
			\begin{tikzpicture}[x=3.65pt,y=3.65pt,xscale=-1]
				\draw[thick,\mycolor,fill=white] (-1,-4) to[out=50,in=-90,looseness=1.4] (-1,-3) to[out=-160,in=0]
				(-2,-3.2) to[out=180,in=-120] (-2.7,-2) to[out=120,in=-120] (-2.9,-1.5)
				to[out=120,in=-120] (-2.95,-1.1) to[out=60,in=-30] (-3.2,-0.9)
				to[out=150,in=-120] (-3,0.4) to[out=60,in=-90] (-3,1.2)
				to[out=90,in=90,looseness=1.8] (3.3,1.2)
				to[out=-90,in=90,looseness=0.8] (2,-2)
				to[out=-90,in=150,looseness=0.8] (3.2,-4) -- cycle;
			\end{tikzpicture}
		};
		\tikzset{declare function={r=0.275;}}
		\path (B) pic[fill=\mycolor,draw=white]{laphuong};
	},
		overlay middle and last ={
		\fill[\mycolor,drop shadow={opacity=0.25,shadow xshift=1pt, shadow yshift=-1.5pt,shadow scale=1.05}] (frame.north west) coordinate (A)--+(2,0)--($(A)+(0,-1.25)$)--cycle;
	}
}
\newenvironment{tomtat}{\begin{kttrongtam}}{\end{kttrongtam}}

%%%==================Môi trường luyện tập=====================%%%
\newcounter{lt}[section] %số điếm của môi trường luyện tập
\newtcolorbox{luyentap}{
	enhanced,breakable,
	before skip=6pt,after skip=6pt,
	boxsep=9pt,
	left=0.75cm,
	right=0.15cm,
	bottom=0pt,
	frame empty,
	frame hidden,
	arc=5pt,
	colback=\mycolor!7,
	sharp corners,
	rounded corners=south,
	colframe=\mycolor!30!blue!45,
	boxrule=0pt,
	borderline north={9pt}{0pt}{\mycolor!50},
	drop fuzzy shadow=gray!75,
	overlay unbroken and first ={
		\fill[\mycolor!80!black,drop shadow={opacity=0.25,shadow xshift=1pt, shadow yshift=-1.5pt,shadow scale=1.05}] (frame.north west) coordinate (A)--+(1.5,0)--($(A)+(0,-1.25)$)--cycle;
		\path ($(A)+(0.5,-0.4)$) coordinate (B) node[font=\LARGE\color{gray!65}]{\faEdit};
		\path ($(B)+(-1pt,1pt)$) node[font=\LARGE\color{white}]{\faEdit};
	},
	code={\refstepcounter{lt}},
	before upper={
		\circlenum{\thelt}
	}
}

%%%=============MT Lưu ý===============%%%%
\usetikzlibrary{shapes.symbols}
\usepackage{paracol}%chia cột có tùy chọn độ rộng
\columnratio{0.7}
\globalcounter{lt}
\newtcolorbox{luuy}{
	enhanced,breakable,
	before skip=6pt,after skip=6pt,
	boxsep=9pt,
	before skip=16pt,
	left=0.75cm,
	right=0.15cm,
	bottom=0pt,
	frame empty,
	frame hidden,
	arc=5pt,
	colback=\mycolor!7,
	sharp corners,
	rounded corners=south,
	colframe=\mycolor!55,
	boxrule=0pt,
	borderline north={9pt}{0pt}{\mycolor!65},
	drop fuzzy shadow=gray!75,
	overlay unbroken and first ={
		\fill[\mycolor!80!black,drop shadow={opacity=0.25,shadow xshift=1pt, shadow yshift=-1.5pt,shadow scale=1.05}] (frame.north west) coordinate (A)--+(1.5,0)--($(A)+(0,-1.25)$)--cycle;
		\draw[\mycolor!80!black,ultra thick,rounded corners=2pt] ($(A)+(6pt,0)$)--++(0,9pt)--++(12pt,0)--+(0,-9pt)
		($(A)+(9pt,0)$)--++(0,5pt)--++(6pt,0)--+(0,-5pt)
		;
		\draw[white,ultra thick,rounded corners=2pt] ($(A)+(18pt,0)$)--++(0,-15pt)--++(-8.5pt,0)--+(0,15pt);
		% \path ($(A)+(0.15,-0.4)$) coordinate (B) node[font=\LARGE\color{gray!65}]{\faPaperclip};
		% \path ($(B)+(-1pt,1pt)$) node[font=\LARGE\color{yellow}]{\faPaperclip};
	}
}

%%%========môi trường note==============%%%
\usepackage{relsize}
\usepackage{amsxtra,amssymb,latexsym,amscd,amsfonts}

\newenvironment{note}{\begin{luuy}%
		{\color{\maunhan}\faBellO\ \relscale{0.9}{\usefont{T5}{put}{b}{n} Lưu ý}.} \ignorespaces}{\end{luuy}}

\newenvironment{nx}{\begin{luuy}%
		{\color{\maunhan}\faPaperclip\ \relscale{0.9}{\usefont{T5}{put}{b}{n} Nhận xét}.} \ignorespaces}{\end{luuy}}

\newenvironment{phuongphap}{\begin{luuy}%
		{\color{\maunhan} \faSpinner\ \relscale{0.9}{\usefont{T5}{put}{b}{n} Phương pháp}.} \ignorespaces}{\end{luuy}}

\newenvironment{phantich}{\begin{luuy}%
		{\color{\maunhan}\faGears\ \relscale{0.9}{\usefont{T5}{put}{b}{n} Phân tích}.} \ignorespaces}{\end{luuy}}


%%%%%%%%%%%%%%%%================Các hộp bổ sung===============================%%%%
%%%%%%%%%%%%%%%%%%%%%%%=======HỘP QANDA============%%%%%%%%%%%%%%%%%%%%%%%%%%%%%%%%%%%
\newcounter{qandacount}[section]
\newtcolorbox{boxqanda}{%
	colback=\mycolor!15,
	colframe=\mycolor!80!black,
	enhanced,
	left=4mm,
	right=4pt,
	top=4mm,
	bottom=6pt,
	after skip=6pt,
	before skip=24pt,
	right skip=3pt,left skip =3pt,
	breakable,
	toprule=0pt,
	bottomrule=0pt,
	rightrule=1.5pt,
	leftrule=1.5pt,
	arc=3pt,
	%	arc is angular,
	sharp corners=northeast,
	drop lifted shadow=\maudam!40!black,
	fontupper=\color{\mycolor!15!black},
	overlay unbroken and first={
		\path 
		([xshift=0.5cm]frame.north west) node[shape=circle,fill=\mycolor!30!orange,inner sep =0pt, outer sep=0pt] (icon)
		{\includegraphics[width=.6cm]{Images/Icons/Icon3dhoidap.png}};
%		\path 
%		([xshift=-3pt]icon.east)	node[anchor=west]{\large\bfseries\sffamily\color{\maunhan}{\theqandacount}}
%		;
	},
	overlay middle and last={
		\path 
		([xshift=0.5cm]frame.north west)  node[shape=circle,fill=\mycolor!30!orange,inner sep =0pt, outer sep=0pt] (icon)
		{\includegraphics[width=.6cm]{Images/Icons/Icon3dhoidap.png}};
%		\path
%		([xshift=-3pt]icon.east)	node[anchor=west]{\large\bfseries\sffamily\color{\maunhan}{\theqandacount}}
%		;
	}	
}

\def\qandabegin{\begin{boxqanda}}
\def\qandaend{\end{boxqanda}}
\newenvironment{hoivadap}{\qandabegin}{\qandaend}
\AtBeginEnvironment{hoivadap}{%
	\refstepcounter{qandacount}
	\setcounter{numTrue}{0}%
	\setcounter{numTrueSol}{0}%
	\renewcommand{\loigiai}[1]{%
		\end{boxqanda} \def\qandaend{} % đưa lời giải ra ngoài hộp
		{\par\noindent\color{\maunhan!50!black}\reflectbox{\Large\WritingHand}\ {\fmmfamily\LARGE Hướng dẫn giải:}}\\%
		#1
	}
}
%%%Môi trường câu hỏi kèm theo%%%
\newcounter{cauhoicount}[section]
\newenvironment{cauhoi}{
\begin{enumerate}[label*=\protect{\small\circlenum[\maunhan]{\arabic*}},wide=0cm,leftmargin=!]}
{\end{enumerate}}
%%%%%%%%%%%%%%%%%%%%%%%=======Hộp Em có biết============%%%%%%%%%%%%%%%%%%%%%%%%%%%%%%%%

\newcounter{ecbcount}[section]
\newtcolorbox{boxecb}{%
	enhanced,
	colback=\mycolor!15,
	left=.8cm,
	right=4pt,
	top=6pt,
	bottom=2pt,
	after skip=6pt,
	before skip=24pt,
	right skip=0cm,
	breakable,
	toprule=1.5pt,
	bottomrule=1.5pt,
	rightrule=1pt,
	leftrule=1pt,
	arc=3pt,
	fonttitle=\Large\bfseries\sffamily\color{white},
	title={Bạn có biết?},
	title after break={Bạn có biết?},
	titlerule style={\mycolor!20!\maunhan,line width=0pt},
	title style={left color=\mycolor!50!\maunhan,right color=\mycolor!75!\maunhan},
	colframe=\mycolor!80!black,
	drop lifted shadow=\maudam!40!black,
	fontupper=\sffamily\color{\mycolor!15!black},
	overlay unbroken and first ={
		\path ($(frame.north west)+(8pt,-5pt)$)  node(Note){\includegraphics[width=1.4cm,angle=5]{Images/Icons/Icon3d_15.png}};
	},
	overlay middle and last={
		\path ($(frame.north west)+(8pt,-5pt)$) node (NT){\includegraphics[width=1.4cm,angle=30]{Images/Icons/Icon3d_15.png}}; 
		
	}
}
\newenvironment{emcobiet}{%
	\refstepcounter{ecbcount}%
	\begin{boxecb}
	}
	{\end{boxecb}}

%%%=============Hộp Ban co biet==========%%%
\newtcolorbox{Bancobiet}[1][\mycolor]{%
	enhanced jigsaw,
	empty,
	breakable,
	vfill before first,
	height fixed for=middle,
	pad at break=1mm,
	before skip =6pt, after skip= 6pt,
	left=6pt,right=6pt,
	boxsep=6pt,
	title= {Em có biết?},
	fonttitle=\color{\mycolor}\large\bfseries\sffamily,
	fontupper=\small\itshape,
	boxed title style ={arc=0pt,outer arc =0pt,frame hidden,colback=\mycolor!20},
	attach boxed title to top left={xshift=1cm,yshift*=0cm-\tcboxedtitleheight/2},
	overlay unbroken and first={%
		interior code={
			%%%=======Xoay 1 góc=======%%%
			\path ($(interior.south west)!0.5!(interior.north east)$) coordinate (G);
			\begin{scope}[transform canvas={rotate around={0:(G)}}]
				\path[fill=#1!25,rounded corners =4pt]
				([shift={(-3pt,0.5pt)}]interior.north west) -- ([shift={(0.5pt,3pt)}]interior.north east)--([shift={(3pt,-0.5pt)}]interior.south east)--([shift={(-0.5pt,-3pt)}]interior.south west)--cycle
				;
			\end{scope}
			\begin{scope}[transform canvas={rotate around={0:(G)}}]
				\path[draw=#1,fill=white!25,rounded corners =4pt,dashed]
				(interior.north west) rectangle (interior.south east);
			\end{scope}
		},
		title code={%
			%%%=====Mảng Tiêu đề=============%%%
			\path  
			(title.north west) coordinate (NW)
			(title.south west) coordinate (SW)
			(title.south east) coordinate (SE)
			(title.north east) coordinate (NE)
			($(NW)!0.5!(SW)$) coordinate (G)
			([xshift=-6pt]G) coordinate (Ga)
			;
			%%%
			\path[fill=#1!90,line cap=round,line join =round]
			 let \p1=($(NW)-(Ga)$),
			 \n1={veclen(\x1,\y1)}
			 in
			(NW)--([xshift=-4pt]NE)
			..controls ++(0:2pt) and ++(90:2pt)..([yshift=-4pt]NE) coordinate (NEt)
			--([yshift=4pt]SE)coordinate (SEt)
			..controls ++(-90:2pt) and ++(0:2pt)..([xshift=-4pt]SE) coordinate (SEt)
			--(SW)--($(Ga)+({-acos(6pt/\n1)}:\n1)$) arc ({-acos(6pt/\n1)}:{acos(6pt/\n1)}:\n1)
			--(NW)
			;
			%%%=====Node Icon=========%%%
			\path[fill=#1!90](Ga) circle (\tcboxedtitleheight/1.88);
			\path (Ga) node{\includegraphics[width=.62cm,trim={0.5cm 0.5cm 0.5cm 0.5cm}]{Images/Icons/Icon_bancobiet.png}};
			%%%Dấu Móc %%%
			\path[fill=#1!90,line cap=round,line join =round](Ga)++(-110:{\tcboxedtitleheight/1.88}) to [out=-120, in= 30] ++(-150:7pt) to [out=40, in= -70]($(Ga)+(-130:{\tcboxedtitleheight/1.88})$)  ;
			%%%=====Node chữ em có biết=========%%%
			\path ($(NW)!0.5!([xshift=4pt]SE)$) node [font=\color{#1!20!white}\large\bfseries\sffamily] {Em có biết?};
		}
		
	},
	overlay middle and last={%
		interior code={
			%%%=======Xoay 1 góc=======%%%
			\path ($(interior.south west)!0.5!(interior.north east)$) coordinate (G);
			\begin{scope}[transform canvas={rotate around={0:(G)}}]
				\path[fill=#1!25,rounded corners =4pt]
				([shift={(-3pt,0.5pt)}]interior.north west) -- ([shift={(0.5pt,3pt)}]interior.north east)--([shift={(3pt,-0.5pt)}]interior.south east)--([shift={(-0.5pt,-3pt)}]interior.south west)--cycle
				;
			\end{scope}
			\begin{scope}[transform canvas={rotate around={0:(G)}}]
				\path[draw=#1,fill=white!25,rounded corners =4pt,dashed]
				(interior.north west) rectangle (interior.south east);
			\end{scope}
		},
		title code={%
			%%%=====Mảng Tiêu đề=============%%%
			\path  
			(title.north west) coordinate (NW)
			(title.south west) coordinate (SW)
			(title.south east) coordinate (SE)
			(title.north east) coordinate (NE)
			($(NW)!0.5!(SW)$) coordinate (G)
			([xshift=-6pt]G) coordinate (Ga)
			;
			%%%
			\path[fill=#1!90,line cap=round,line join =round]
			let \p1=($(NW)-(Ga)$),
			\n1={veclen(\x1,\y1)}
			in
			(NW)--([xshift=-4pt]NE)
			..controls ++(0:2pt) and ++(90:2pt)..([yshift=-4pt]NE) coordinate (NEt)
			--([yshift=4pt]SE)coordinate (SEt)
			..controls ++(-90:2pt) and ++(0:2pt)..([xshift=-4pt]SE) coordinate (SEt)
			--(SW)--($(Ga)+({-acos(6pt/\n1)}:\n1)$) arc ({-acos(6pt/\n1)}:{acos(6pt/\n1)}:\n1)
			--(NW)
			;
			%%%=====Node Icon=========%%%
			\path[fill=#1!90](Ga) circle (\tcboxedtitleheight/1.88);
			\path (Ga) node{\includegraphics[width=.62cm,trim={0.5cm 0.5cm 0.5cm 0.5cm}]{Images/Icons/Icon_bancobiet.png}};
			%%%Dấu Móc %%%
			\path[fill=#1!90,line cap=round,line join =round](Ga)++(-110:{\tcboxedtitleheight/1.88}) to [out=-120, in= 30] ++(-150:7pt) to [out=40, in= -70]($(Ga)+(-130:{\tcboxedtitleheight/1.88})$)  ;
			%%%=====Node chữ em có biết=========%%%
			\path ($(NW)!0.5!([xshift=4pt]SE)$) node [font=\color{#1!20!white}\large\bfseries\sffamily] {Em có biết?};
		}
	}%
}

%%%%%%%%%%%%%%%%%%%%%%=======Hộp tổng kết============%%%%%%%%%%%%%%%%%%%%%%%%%%%%%%%%

\newcounter{tkcount}[section]
\newtcolorbox[auto counter]{boxtk}{%
	empty,
	colback=\mycolor!15,
	enhanced,
	left=3.5mm,
	right=6pt,
	top=6pt,
	bottom=6pt,
	after skip=6pt,
	before skip=6pt,
	right skip=0cm,
	breakable,
	toprule=0pt,
	bottomrule=1.5pt,
	rightrule=1.5pt,
	leftrule=0pt,
	arc=.3cm,
	sharp corners=northwest,
	sharp corners=southwest,
	colframe=\maunhan!90!white,
	drop lifted shadow=\maudam!40!black,
	fontupper=\sffamily\color{\mycolor!15!black},
	overlay unbroken and first ={
		\path[fill=violet!80!orange] ($(frame.north west)+(0.5cm,-6pt)$) coordinate (tron)  circle (.4cm);
		\path (tron)  node[inner sep=0pt](Note){\includegraphics[width=0.7cm]{Images/Icons/Icon3d_emdahoc.png}};
	},
	overlay middle and last={
		\path[fill=violet!80!orange] ($(frame.north west)+(0.5cm,-6pt)$) coordinate (tron)  circle (.4cm);
		\path (tron)  node[inner sep=0pt](Note){\includegraphics[width=0.7cm]{Images/Icons/Icon3d_emdahoc.png}};
	}
}

\newenvironment{tongket}[1]{%
	\refstepcounter{tkcount}	
	\begin{boxtk}
		{\hspace*{5mm}\large\bfseries\fontfamily{qag}\selectfont\color{\maunhan}{#1}}}
	{\end{boxtk}}
%%%%%%%%%%%%%%%%%==============Hộp Khởi động========================%%%%%%%%%%%%%%%%
\newcounter{kdcount}[section]
\newtcolorbox{boxkd}{%
	enhanced,
	frame empty,
	colback=\maunhan!15,
	left=4mm,
	right=6pt,
	top=5mm,
	bottom=6pt,
	after skip=6pt,
	before skip=24pt,
	right skip=6pt,
	arc=9pt,
	arc is angular,
	sharp corners =northeast,sharp corners =northwest,sharp corners =southwest,
	breakable,
%	toprule=0pt,
%	bottomrule=1.5pt,
%	rightrule=1.5pt,
%	leftrule=1.5pt,
%	colframe=\maunhan,
	drop lifted shadow=\maudam!40!black,
	fontupper=\sffamily\color{\maunhan!40!black},	
	overlay unbroken and first={
		\path 
		([shift={(0.7cm,-5pt)}]frame.north west) node[shape=circle,fill=\maunhan!60!violet!90,inner sep=0pt, outer sep=0pt] (note)
		{\includegraphics[width=0.7cm,trim={0.3cm 0.3cm 0.3cm 0.3cm}]{Images/Icons/Icon3d_16.png} }
		;
	},
	overlay middle and last=={
		\path 
		([shift={(0.7cm,-5pt)}]frame.north west) node[shape=circle,fill=\maunhan!60!violet!90,inner sep=0pt, outer sep=0pt] (note)
		{\includegraphics[width=0.7cm,trim={0.3cm 0.3cm 0.3cm 0.3cm}]{Images/Icons/Icon3d_16.png} }
		;
	}
}
\def\kdbegin{\begin{boxkd}}
	\def\kdend{\end{boxkd}}
\newenvironment{kd}
{\kdbegin}
{\kdend}
\AtBeginEnvironment{kd}{%
	\refstepcounter{kdcount}
	\setcounter{numTrue}{0}%
	\setcounter{numTrueSol}{0}%
	\renewcommand{\loigiai}[1]{%
		\par\noindent
		\end{boxkd} \def\kdend{} % đưa lời giải ra ngoài hộp
		{\par\noindent\color{\mauphu}\reflectbox{\Large\WritingHand}\ {\fmmfamily\LARGE Lời giải:}} 
		\par\noindent
		#1
	}
}			
%%%%%%%%%%%%%%%%%==============Hộp lý thuyết========================%%%%%%%%%%%%%%%%
\newcounter{ltcount}[section]
\newtcolorbox{boxlt}{%
	enhanced,
	frame empty,
	colback=\mycolor!15,
	left=3mm,
	right=2mm,
	top=6mm,
	bottom=6pt,
	after skip=5mm,
	before skip=5mm,
	right skip=0cm,
	breakable,
	toprule at break=\tcboxedtitleheight,
	toprule=0pt,
	arc=3pt,
	sharp corners=northwest,sharp corners=northeast,
	bottomrule=1pt,
	rightrule=1pt,
	leftrule=1pt,
	colframe=\mycolor,
	drop lifted shadow=\maudam!40!black,
	overlay unbroken and first={
		\path 
		([shift={(0.6cm,-0.1cm)}]frame.north west) node[fill=violet!80,inner sep=0pt,outer sep =-1pt,shape=circle] (note)
		{\includegraphics[width=.62cm,trim={0.5cm 0.5cm 0.5cm 0.5cm}]{Images/Icons/light.pdf}};
		\draw[\maunhan!90!black,rounded corners=4pt,ultra thick] (frame.north west)--(frame.south west)-- (frame.south east)--(frame.north east);
	},
	overlay middle and last ={
		\path 
		([shift={(0.6cm,-0.1cm)}]frame.north west) node[fill=violet!80,inner sep=0pt,outer sep =-1pt,shape=circle] (note)
		{\includegraphics[width=.62cm,trim={0.5cm 0.5cm 0.5cm 0.5cm}]{Images/Icons/light.pdf}};
		\draw[\maunhan!90!black,rounded corners=4pt,ultra thick] (frame.north west)--(frame.south west)-- (frame.south east)--(frame.north east); 
	}
}
\newenvironment{ghinho}{\begin{boxlt}}{\end{boxlt}}
%%%=============Icon Bóng đèn ===========
\NewDocumentCommand{\bongden}{O{0.7}O{violet!80}}{
	\begin{tikzpicture}[declare function={d=2.0;h=0.45*d;k=#1;R=1.2 cm ;w=R*sin(20);},line join=round,line cap=round,scale=k]
		\path (0,0) circle (2pt) coordinate (O)
		($(O)+(-d,0)$) circle (2pt) coordinate (Aa)
		($(O)+(d,0)$) circle (2pt) coordinate (Ab)
		($(O)+(0,-1.3*h)$) circle (2pt) coordinate (B)
		($(B)+(-0.75*d,0)$) circle (2pt) coordinate  (Ba)
		($(B)+(0.75*d,0)$) circle (2pt) coordinate (Bb)
		($(O)+(0,-2.3*h)$) circle (2pt) coordinate (C)
		($(C)+(-0.45*d,0)$) circle (2pt) coordinate  (Ca)
		($(C)+(0.45*d,0)$) circle (2pt) coordinate (Cb)
		($(O)+(0,-2.4*h)$) circle (2pt) coordinate (D)
		($(D)+(-0.5*d,0)$) circle (2pt) coordinate  (Da)
		($(D)+(0.5*d,0)$) circle (2pt) coordinate (Db)
		($(Ca)!0.5!(Cb)$) coordinate (Co)
		;
		%%%=====Vẽ bóng tròn=========%%%
		\path[fill=orange!70!white,even odd rule]
		(O) circle (R)
		($(O)+(250:R)$) --([xshift=-w]Co)--(Ca)..controls ++(90:0.8cm) and ++(-60:0.05cm)..
		(Ba)..controls ++(120:0.01cm) and ++(-90:1.4cm)..(Aa)
		--(Aa) arc (180:0:d) coordinate (Ab)
		..controls ++(-90:1.4cm) and ++(60:0.01cm)..(Bb)
		..controls ++(-120:0.05cm) and ++(90:0.8cm)..(Cb)--([xshift=w]Co)--($(O)+ (290:R)$)--cycle
		;
		%%%=====Vẽ tia sáng=========%%%
		\foreach \g/\p in {-30/west,0/west,30/west,60/west,90/west,120/east,150/east,180/east,210/east}{
			\path (O)--++(\g:1.05*d) node[anchor=\p,sloped,midway,pos=1,scale=k]{%
				\tikz{%
					\path[fill=yellow!20!orange,rounded corners=3pt]
					(0,0)--++(0.2*d,0)--++(-90:8pt)--++(180:0.2*d)--cycle;
				}
			};
		}
		%%%=====Vẽ chuôi bóng hình chữ nhật=========%%%
		\path[fill=white!30!gray,rounded corners =8pt] (Da)--++(-90:40pt)  coordinate (Dc)--++(0:d)  coordinate (Dd)--(Db);
		%%%=====Vẽ chuôi bóng tròn=========%%%
		\path[fill=white!90!gray] ([xshift=8pt]Dc)..controls ++(-90:2.2cm) and ++(-90:2.2cm)..([xshift=-8pt]Dd);
		\path[fill=orange!30!white] (O) circle ({R-0.7cm}) ;
		\path[fill=#2] ([yshift=3pt]$(O)+(250:R)$) rectangle ([yshift=-3pt]$(O)+(290:R)$);
	\end{tikzpicture}
}

%%%=============Hộp vi du don gian==========%%%
\newtcolorbox{hopvidu}[1][\mycolor]{%
	enhanced jigsaw,
	opacityback=0,
	empty,
	breakable,
	before skip =6pt, after skip= 6pt,
	left=6pt,right=6pt,
	boxsep=6pt,
	fontupper=\small\itshape,
	overlay unbroken and first={%
		interior code={
			%%%=======Xoay 1 góc=======%%%
			\path ($(interior.south west)!0.5!(interior.north east)$) coordinate (G);
			\begin{scope}[transform canvas={rotate around={0:(G)}}]
				\path[fill=#1!45,rounded corners =4pt]
				([shift={(-3pt,0.5pt)}]interior.north west) -- ([shift={(0.5pt,3pt)}]interior.north east)--([shift={(3pt,-0.5pt)}]interior.south east)--([shift={(-0.5pt,-3pt)}]interior.south west)--cycle
				;
			\end{scope}
			\begin{scope}[transform canvas={rotate around={0:(G)}}]
				\path[draw=#1,fill=#1!10!white!25,rounded corners =4pt,dashed]
				(interior.north west) rectangle (interior.south east);
			\end{scope}
		}	
	},
	overlay middle and last={%
		interior code={
			%%%=======Xoay 1 góc=======%%%
			\path ($(interior.south west)!0.5!(interior.north east)$) coordinate (G);
			\begin{scope}[transform canvas={rotate around={0:(G)}}]
				\path[fill=#1!25,rounded corners =4pt]
				([shift={(-3pt,0.5pt)}]interior.north west) -- ([shift={(0.5pt,3pt)}]interior.north east)--([shift={(3pt,-0.5pt)}]interior.south east)--([shift={(-0.5pt,-3pt)}]interior.south west)--cycle
				;
			\end{scope}
			\begin{scope}[transform canvas={rotate around={0:(G)}}]
				\path[draw=#1,fill=white!25,rounded corners =4pt,dashed]
				(interior.north west) rectangle (interior.south east);
			\end{scope}
		}
	}%
}
\newtcolorbox{hopdongian}[1][\mycolor]{%
	enhanced jigsaw,
	frame empty,
	colback=#1!20,
	arc is angular,
	arc=4pt
	}

%%Lệnh hình phải
%\NewDocumentCommand{\Rightpicture}{O{0.45}O{4pt}mm}{%
%	\adjustbox{valign=t}{
%		\begin{minipage}[htp!]{#1\linewidth}
%			\vspace*{#2}
%			#3
%		\end{minipage}
%	}
%	\hfill
%	\adjustbox{valign=t}{
%		\begin{minipage}[htp!]{0.98\linewidth-#1\linewidth}
%			#4
%		\end{minipage}
%	}
%}
%
%\newenvironment{vidu}[1][\mycolor]{
%	\begin{hopvidu}[#1]
%		\par\noindent\indam[#1]{Ví dụ:}
%	}{\end{hopvidu}}
%\AtBeginEnvironment{ex}{\renewcommand{\dotEX}{}}




